\section{The Reduction and Its Missing Link}
\label{sec:verified-reduction}

\Cref{sec:simulation-target} isolates the local change behind the LMSS
reduction: replace \(\mathsf{ODec}_{\mathsf{real}}\) by the
plaintext-centered \(\mathsf{ODec}_{\mathsf{sim}}\).
The latter can be implemented without the secret key and hence inside an
IND-CPA reduction.
For one call, \Cref{eq:epsilon-nf} bounds the KL cost.
We now lift that local comparison to the complete adaptive experiment.

\subsection{Three games}

Fix an IND-CPAD adversary \(\cA\) making at most \(q\) decryption queries.
Consider the following games:
\begin{align*}
 G_0 &=
   \mathsf{IND\text{-}CPAD}_{\mathsf{NF}_{\gamma}[S],q}(\cA),\\
 G_1 &=
   \mathsf{IND\text{-}CPAD}^{\mathsf{sim}}_
     {\mathsf{NF}_{\gamma}[S],q}(\cA),\\
 G_2 &=
   \mathsf{IND\text{-}CPA}_{S}(\cB_{\cA,q}).
\end{align*}
Game \(G_0\) is the real experiment.
Game \(G_1\) replaces the real decryption oracle by
\(\mathsf{ODec}_{\mathsf{sim}}\) from \Cref{eq:real-sim-decrypt}.
This simulated-decryption game is the intermediate game used by
LMSS~\cite{lmss}.
Game \(G_2\) runs the concrete reduction from
\Cref{alg:ind-cpa-reduction}.

The second transition is exact.
Once decryption answers are centered at the common recorded plaintext, they no
longer require the secret key or the hidden challenge bit.
The reduction can therefore forward encryption queries to its IND-CPA oracle,
perform evaluation and table bookkeeping locally, and simulate every permitted
decryption query with the same distribution as \(G_1\).
Consequently,
\[
 \WinProb[G_1]=\WinProb[G_2].
\]

The entire cryptographic proof is thus the chain
\begin{equation}
\label{eq:conceptual-game-chain}
 G_0
 \xrightarrow{\sqrt{q\epsilon_{\mathsf{nf}}/2}}
 G_1
 \xrightarrow{0}
 G_2.
\end{equation}
The first arrow is the only paid hop.
Together with the assumed IND-CPA bound for \(G_2\), it immediately implies
\Cref{thm:main}.

\subsection{The tempting one-line proof}
\label{sec:paid-hop}

\Cref{eq:epsilon-nf} bounds the conditional KL cost of every reachable
decryption query whose two recorded plaintexts agree by
\(\epsilon_{\mathsf{nf}}\).
The two games agree when those plaintexts differ and when an assertion fails,
so those branches spend no budget.

It is tempting to finish in one sentence:
apply the KL chain rule over the \(q\) decryption interactions, obtain total KL cost
\(q\epsilon_{\mathsf{nf}}\), and apply Pinsker once.
This gives
\begin{equation}
\label{eq:central-game-hop}
 \TVD(G_0,G_1)
 \leq \sqrt{q\epsilon_{\mathsf{nf}}/2}.
\end{equation}
This is the KL-composition argument used by LMSS~\cite{lmss}.
A machine-checked proof must identify the relevant conditional distributions
and establish the one-call bound after every adversarially chosen history.

\subsection{Where adaptivity enters}
\label{sec:adaptive-invariant}

The adversary's decryption calls do not form a fixed product experiment.
After seeing an answer, \(\cA\) may change its state, choose a different oracle,
or decide whether and where to make its next decryption query.
In SSProve, \(\cA\) is one stateful program with calls embedded in its control
flow, not a list of \(q\) pre-existing query phases.
Thus ``the first \(q\) decryption calls'' is a semantic notion, and the
conditional KL bound must hold after every reachable history.

The proof also depends on a state invariant \(\Phi\) relating the real and
simulated executions.
It records initialized, good keys; agreement of their public data and hidden
challenge bits; validity of the challenge table under encryption and
evaluation; and agreement of the decryption counters.
In particular, for every reachable row
\((m_0,m_1,c)\) with \(c=\mathsf{Some}(\bar c,e)\), the plaintext on the
selected branch \(b\) satisfies
\[
 d(\mathsf{dec}_0(sk,c),m_b)\leq e.
\]
Initialization establishes \(\Phi\), while encryption, evaluation, and the
deterministic prefix of decryption preserve it.
This is what makes the one-call KL calculation available after an arbitrary
adaptive history.

A machine-checked proof of \Cref{eq:central-game-hop} therefore needs a bridge
with three properties:
\begin{enumerate}
\item it retains conditional KL costs instead of converting each call to
statistical distance;
\item it exposes successive calls of an arbitrary adaptive oracle program; and
\item it threads \(\Phi\) through the exposed calls.
\end{enumerate}
\Cref{sec:pyth-toolkit} develops the Pythagorean judgment and invariant rules.
\Cref{sec:trace-compiler} verifies the selected-call compiler and proves the
generic local-to-adaptive theorem independently of FHE.  The instantiation in
\Cref{sec:closing-hop} combines them to derive the paid arrow in
\Cref{eq:conceptual-game-chain}.
